% ---------------------------------------------------------------------------
% Front matter: cover, abstract, contents
% ---------------------------------------------------------------------------
\begin{titlepage}
\thispagestyle{empty}
\begin{tikzpicture}[remember picture,overlay]
  \fill[paper] (current page.north west) rectangle (current page.south east);
  \fill[deepgreen] (current page.north west) rectangle ([yshift=-8mm]current page.north east);
  \fill[accent] ([yshift=-8mm]current page.north west) rectangle ([yshift=-9mm]current page.north east);
  % faint pixel rule at the foot
  \fill[accentlight!55] ([yshift=16mm]current page.south west) rectangle ([yshift=16.4mm]current page.south east);
\end{tikzpicture}

\vspace*{12mm}
{\sffamily\scriptsize\color{muted}\MakeUppercase{Hopper Protocol}\hfill Technical Whitepaper}
\vspace{2mm}

\vspace{14mm}
\begin{center}
  \includegraphics[width=46mm]{figures/hopper-logo.png}\\[10mm]
  {\sffamily\bfseries\color{deepgreen}\fontsize{60}{64}\selectfont HOPPER}\\[3mm]
  {\sffamily\Large\color{ink} A Satoshi-Native Game Protocol\\[1.2mm] on Bitcoin}\\[9mm]
  {\color{accent}\rule{54mm}{1.1pt}}\\[7mm]
  {\large\itshape One satoshi. One life.\\[1mm] The chain is the referee.}
\end{center}

\vfill
\begin{center}
  \begin{tabular}{@{}c@{}}
    {\sffamily\small\color{deepgreen}\textbf{Wire version} \code{0x08}}\\[1.2mm]
    {\sffamily\small\color{muted}Specification 3.0 \textperiodcentered{} October 2026}\\[1.2mm]
    {\sffamily\small\color{muted}The HOPPER contributors}
  \end{tabular}
\end{center}
\vspace*{10mm}
\end{titlepage}

% ---------------------------------------------------------------------------
% Abstract
% ---------------------------------------------------------------------------
\pagenumbering{roman}
\thispagestyle{plain}
\section*{Abstract}
\addcontentsline{toc}{section}{Abstract}

\begin{notebox}
\noindent
HOPPER is a game played by a satoshi. A frog is an ordinal sat whose public
life is written into Bitcoin itself: a mint inscription creates it, each
crossing publishes a 55-byte state inscription on the same sat, and the
direction carried in that inscription is judged against the one direction
the next block wraps to. Custody is enforced by a four-leaf Taproot output
with a nothing-up-my-sleeve internal key. Three crossing leaves, one per
direction, each check the owner's signature and carry the protocol envelope;
the fourth leaf is an ordinary owner signature for transfers. The v8 fold
verifies every crossing body against the raw transaction that published it;
if a proof is missing, the art says so rather than guessing. The road's
difficulty runs in generations: it rises from five crossings to a ceiling of
twenty-five and then resets. One released v7 Hopper is grandfathered for
custody and play. No server signs, no external origin is needed to render the
frog, and the indexer is a cache, not an authority.

This paper specifies the wire format, the generation schedule, the
block-wide wrap, the tapscript tree and transaction shapes that play the
game, and the standalone indexer and JavaScript wallet that implement them.
It is an exposition of the normative specification in \code{SPEC.md}; where
the two documents differ, \code{SPEC.md} governs.
\end{notebox}

\vspace{4mm}
\noindent\textbf{\sffamily Keywords:} Bitcoin, ordinals, inscriptions,
Taproot, tapscript, ordinal theory, self-custody, PWA.

\vspace{2mm}
\noindent\textbf{\sffamily Normative source.} The repository specification
\code{hopper-protocol/SPEC.md} is the single normative source for the
consensus rules, byte layouts, constants, API surface and configuration
described here. Code snippets are illustrative; byte strings, constants and
acceptance conditions are quoted from that document.

\vspace{6mm}
{\small\color{muted}\noindent
The historical mint magic \code{FRG} and the JSON keys
\code{frog}/\code{frogger} are retained because they are already committed to
Bitcoin. The current state magic is \code{FRG8CROSS}; the legacy
\code{FRG7CROSS} is accepted only for the one grandfathered Hopper. The
interface name is \textsf{HOPPER}; the crossing action is \code{HOP} in the
interface and \code{cross} on the wire.}

\clearpage

% ---------------------------------------------------------------------------
% Contents
% ---------------------------------------------------------------------------
\thispagestyle{plain}
\renewcommand{\contentsname}{\sffamily\color{deepgreen}Contents}
\tableofcontents
\vfill
\begin{notebox}
\small
\textbf{Reading guide.} Sections~\ref{sec:parameters}--\ref{sec:custody}
are the consensus core. Sections~\ref{sec:indexer} and~\ref{sec:wallet}
describe infrastructure that can be replaced without changing the protocol.
Section~\ref{sec:security} states what the construction does and does not
protect against.
\end{notebox}
\clearpage
